% !Mode:: "TeX:UTF-8"
% 问题重述草稿：依据题目正文与数据说明整理。
% 本文件为章节片段，可由 paper/main.tex 使用 % !Mode:: "TeX:UTF-8"
% 问题重述草稿：依据题目正文与数据说明整理。
% 本文件为章节片段，可由 paper/main.tex 使用 % !Mode:: "TeX:UTF-8"
% 问题重述草稿：依据题目正文与数据说明整理。
% 本文件为章节片段，可由 paper/main.tex 使用 % !Mode:: "TeX:UTF-8"
% 问题重述草稿：依据题目正文与数据说明整理。
% 本文件为章节片段，可由 paper/main.tex 使用 \input{sections/drafts/problem-restatement.tex} 引入。
% 仅包含问题背景与问题提出，不包含具体模型、计算结果或结论。

\section{问题重述}
\label{sec:problem-restatement}

\subsection{问题背景}

近年来，大语言模型在自然语言处理、科学研究和智能应用等领域快速发展，经证明发现模型能力通常受到参数规模、训练数据量和训练算力的共同影响。其中，标度律是用以刻画模型损失随规模变化的经验关系，而经典标度律则是主要用于描述参数量 $N$、训练数据量 $D$ 与验证集交叉熵损失 Loss 之间的关系。但是随着大语言模型的长期实践发现，模型性能并不只取决于规模因素，训练数据质量、不同领域数据的配比、上下文长度以及模型架构等因素同样会产生重要影响。特别是在高质量语料日益稀缺、数据污染和“数据墙”问题逐渐突出的背景下，仅依靠扩大模型规模或增加训练数据，已经难以全面解释模型能力的变化。

但是与此同时，增加参数量、扩充训练数据、提升数据质量以及处理更长上下文都会带来额外的算力消耗；故而在算力预算有限的情况下，模型规模、数据规模、数据质量和领域配比之间存在资源竞争与相互替代关系。因此，本题要求综合利用附件中的数据质量信号、配方实验、模型训练日志、历史模型元数据和多维 Benchmark 评测数据，建立“数据质量与领域配比—模型 Loss—模型能力—算力配置”之间的数学联系。四个问题依次研究数据质量和领域配比、广义标度律、算力约束下的资源优化以及模型能力演进预测，前一问题的分析结果将作为后一问题的重要基础。

\subsection{问题提出}

\subsubsection{问题一：数据质量评价、质量冲突消解与领域配比建模}

问题一主要研究训练数据质量及领域配比对模型性能的影响。利用附件 $A_1$--$A_3$ 中的 22 个质量指标，需要对不同方向的指标进行预处理、标准化和正负向统一，构造样本级、语料级或领域级质量评分 $Q$，并依据 $A_{16}$ 建立质量信号领域与 17 个训练领域之间的对应关系。同时，需要识别同一文本上不同质量指标评价明显不一致的情况，定义质量冲突并分析其产生原因，在抽样集和扩展集上检验冲突规律的稳定性。在此基础上，利用 $A_4$--$A_{15}$ 中的 17 维领域配比和交叉熵损失数据，建立领域配比向量 $\boldsymbol{p}$ 与 Loss 之间的定量关系，其中
\begin{equation}
  \boldsymbol{p}=(p_1,p_2,\ldots,p_{17}),\qquad
  p_i\geq 0,\qquad \sum_{i=1}^{17}p_i=1.
  \label{eq:problem-restatement-simplex}
\end{equation}
需要分析不同领域及其组合对模型性能的影响，并通过检验数据和外推数据验证模型的可靠性与稳健性。

\subsubsection{问题二：跨维度数据融合与广义标度律}

问题二要求在经典标度律的基础上，将问题一得到的数据质量 $Q$ 和领域配比 $\boldsymbol{p}$ 纳入模型，建立可以同时考虑参数量 $N$、训练数据量 $D$、数据质量和领域组成的广义标度律。模型拟合主要使用附件 $B_1$ 中的真实训练轨迹，并结合 $B_2$、$B_3$ 进行轨迹或模型外部验证，利用 $B_4$、$B_5$ 检验模型在跨模型和文献数据上的泛化能力，再通过 $B_6$--$B_8$ 的质量补充实验分析数据质量的作用，利用 $B_9$、$B_{10}$ 讨论大规模模型的外推边界。在模型估计和验证的基础上，还需要比较参数规模、数据规模、数据质量和领域配比的边际效用与弹性，分析不同领域之间的替代或互补关系，并推导质量提升与参数规模扩张之间的可计算等价条件，同时明确半合成数据、插值数据和外推数据的可信度边界。

\subsubsection{问题三：算力约束下的多维资源联合优化与结构性转移}

问题三研究有限算力预算下的资源配置问题，需要在参数量 $N$、训练数据量 $D$、数据质量 $Q$ 和领域配比 $\boldsymbol{p}$ 之间进行联合分配，使模型预测 Loss 尽可能低。优化模型应同时考虑基础训练成本、数据质量提升成本和长上下文注意力成本，并比较指数型、幂函数型和对数渐进型质量成本函数对最优解的影响。上下文长度 $L_{\mathrm{ctx}}$ 由模型架构和任务需求外生给定，需要依据附件 $C_7$ 确定其合理范围并进行敏感性分析，同时设置低、中、高不同量级的算力预算，考察预算变化时各类资源的最优分配比例是否发生明显改变，进一步给出结构性转移的判定标准以及训练成本与注意力成本相当的临界条件。

\subsubsection{问题四：技术演进分析与模型能力前沿预测}

问题四利用附件 $C$ 中的模型规模、训练算力、训练数据量、发布时间、开源状态和多维 Benchmark 数据，分析大语言模型能力提升的来源，并预测未来能力的前沿边界。首先，需要将能力提升分解为参数规模、数据规模和算力增长带来的规模扩张效应，以及架构改进、数据工程和对齐方法等带来的非规模技术进步，并分别量化和比较两类因素的贡献。其次，需要明确综合 Benchmark 指标、开源模型筛选标准、\emph{pretrained} 与 \emph{chat/finetuned} 模型的分类方式以及时间轴口径，并利用 $C_5$ 或 $C_6$ 建立 Loss 与 Benchmark 之间的桥接关系，分析桥接误差对结论的影响。最后，结合 $C_1$--$C_4$ 的总体信息和 $C_8$ 的逐任务评测结果，在算力增长放缓的情景下预测未来 12 个月或 24 个月开源大语言模型的能力前沿边界，并给出相应的不确定性分析。
 引入。
% 仅包含问题背景与问题提出，不包含具体模型、计算结果或结论。

\section{问题重述}
\label{sec:problem-restatement}

\subsection{问题背景}

近年来，大语言模型在自然语言处理、科学研究和智能应用等领域快速发展，经证明发现模型能力通常受到参数规模、训练数据量和训练算力的共同影响。其中，标度律是用以刻画模型损失随规模变化的经验关系，而经典标度律则是主要用于描述参数量 $N$、训练数据量 $D$ 与验证集交叉熵损失 Loss 之间的关系。但是随着大语言模型的长期实践发现，模型性能并不只取决于规模因素，训练数据质量、不同领域数据的配比、上下文长度以及模型架构等因素同样会产生重要影响。特别是在高质量语料日益稀缺、数据污染和“数据墙”问题逐渐突出的背景下，仅依靠扩大模型规模或增加训练数据，已经难以全面解释模型能力的变化。

但是与此同时，增加参数量、扩充训练数据、提升数据质量以及处理更长上下文都会带来额外的算力消耗；故而在算力预算有限的情况下，模型规模、数据规模、数据质量和领域配比之间存在资源竞争与相互替代关系。因此，本题要求综合利用附件中的数据质量信号、配方实验、模型训练日志、历史模型元数据和多维 Benchmark 评测数据，建立“数据质量与领域配比—模型 Loss—模型能力—算力配置”之间的数学联系。四个问题依次研究数据质量和领域配比、广义标度律、算力约束下的资源优化以及模型能力演进预测，前一问题的分析结果将作为后一问题的重要基础。

\subsection{问题提出}

\subsubsection{问题一：数据质量评价、质量冲突消解与领域配比建模}

问题一主要研究训练数据质量及领域配比对模型性能的影响。利用附件 $A_1$--$A_3$ 中的 22 个质量指标，需要对不同方向的指标进行预处理、标准化和正负向统一，构造样本级、语料级或领域级质量评分 $Q$，并依据 $A_{16}$ 建立质量信号领域与 17 个训练领域之间的对应关系。同时，需要识别同一文本上不同质量指标评价明显不一致的情况，定义质量冲突并分析其产生原因，在抽样集和扩展集上检验冲突规律的稳定性。在此基础上，利用 $A_4$--$A_{15}$ 中的 17 维领域配比和交叉熵损失数据，建立领域配比向量 $\boldsymbol{p}$ 与 Loss 之间的定量关系，其中
\begin{equation}
  \boldsymbol{p}=(p_1,p_2,\ldots,p_{17}),\qquad
  p_i\geq 0,\qquad \sum_{i=1}^{17}p_i=1.
  \label{eq:problem-restatement-simplex}
\end{equation}
需要分析不同领域及其组合对模型性能的影响，并通过检验数据和外推数据验证模型的可靠性与稳健性。

\subsubsection{问题二：跨维度数据融合与广义标度律}

问题二要求在经典标度律的基础上，将问题一得到的数据质量 $Q$ 和领域配比 $\boldsymbol{p}$ 纳入模型，建立可以同时考虑参数量 $N$、训练数据量 $D$、数据质量和领域组成的广义标度律。模型拟合主要使用附件 $B_1$ 中的真实训练轨迹，并结合 $B_2$、$B_3$ 进行轨迹或模型外部验证，利用 $B_4$、$B_5$ 检验模型在跨模型和文献数据上的泛化能力，再通过 $B_6$--$B_8$ 的质量补充实验分析数据质量的作用，利用 $B_9$、$B_{10}$ 讨论大规模模型的外推边界。在模型估计和验证的基础上，还需要比较参数规模、数据规模、数据质量和领域配比的边际效用与弹性，分析不同领域之间的替代或互补关系，并推导质量提升与参数规模扩张之间的可计算等价条件，同时明确半合成数据、插值数据和外推数据的可信度边界。

\subsubsection{问题三：算力约束下的多维资源联合优化与结构性转移}

问题三研究有限算力预算下的资源配置问题，需要在参数量 $N$、训练数据量 $D$、数据质量 $Q$ 和领域配比 $\boldsymbol{p}$ 之间进行联合分配，使模型预测 Loss 尽可能低。优化模型应同时考虑基础训练成本、数据质量提升成本和长上下文注意力成本，并比较指数型、幂函数型和对数渐进型质量成本函数对最优解的影响。上下文长度 $L_{\mathrm{ctx}}$ 由模型架构和任务需求外生给定，需要依据附件 $C_7$ 确定其合理范围并进行敏感性分析，同时设置低、中、高不同量级的算力预算，考察预算变化时各类资源的最优分配比例是否发生明显改变，进一步给出结构性转移的判定标准以及训练成本与注意力成本相当的临界条件。

\subsubsection{问题四：技术演进分析与模型能力前沿预测}

问题四利用附件 $C$ 中的模型规模、训练算力、训练数据量、发布时间、开源状态和多维 Benchmark 数据，分析大语言模型能力提升的来源，并预测未来能力的前沿边界。首先，需要将能力提升分解为参数规模、数据规模和算力增长带来的规模扩张效应，以及架构改进、数据工程和对齐方法等带来的非规模技术进步，并分别量化和比较两类因素的贡献。其次，需要明确综合 Benchmark 指标、开源模型筛选标准、\emph{pretrained} 与 \emph{chat/finetuned} 模型的分类方式以及时间轴口径，并利用 $C_5$ 或 $C_6$ 建立 Loss 与 Benchmark 之间的桥接关系，分析桥接误差对结论的影响。最后，结合 $C_1$--$C_4$ 的总体信息和 $C_8$ 的逐任务评测结果，在算力增长放缓的情景下预测未来 12 个月或 24 个月开源大语言模型的能力前沿边界，并给出相应的不确定性分析。
 引入。
% 仅包含问题背景与问题提出，不包含具体模型、计算结果或结论。

\section{问题重述}
\label{sec:problem-restatement}

\subsection{问题背景}

近年来，大语言模型在自然语言处理、科学研究和智能应用等领域快速发展，经证明发现模型能力通常受到参数规模、训练数据量和训练算力的共同影响。其中，标度律是用以刻画模型损失随规模变化的经验关系，而经典标度律则是主要用于描述参数量 $N$、训练数据量 $D$ 与验证集交叉熵损失 Loss 之间的关系。但是随着大语言模型的长期实践发现，模型性能并不只取决于规模因素，训练数据质量、不同领域数据的配比、上下文长度以及模型架构等因素同样会产生重要影响。特别是在高质量语料日益稀缺、数据污染和“数据墙”问题逐渐突出的背景下，仅依靠扩大模型规模或增加训练数据，已经难以全面解释模型能力的变化。

但是与此同时，增加参数量、扩充训练数据、提升数据质量以及处理更长上下文都会带来额外的算力消耗；故而在算力预算有限的情况下，模型规模、数据规模、数据质量和领域配比之间存在资源竞争与相互替代关系。因此，本题要求综合利用附件中的数据质量信号、配方实验、模型训练日志、历史模型元数据和多维 Benchmark 评测数据，建立“数据质量与领域配比—模型 Loss—模型能力—算力配置”之间的数学联系。四个问题依次研究数据质量和领域配比、广义标度律、算力约束下的资源优化以及模型能力演进预测，前一问题的分析结果将作为后一问题的重要基础。

\subsection{问题提出}

\subsubsection{问题一：数据质量评价、质量冲突消解与领域配比建模}

问题一主要研究训练数据质量及领域配比对模型性能的影响。利用附件 $A_1$--$A_3$ 中的 22 个质量指标，需要对不同方向的指标进行预处理、标准化和正负向统一，构造样本级、语料级或领域级质量评分 $Q$，并依据 $A_{16}$ 建立质量信号领域与 17 个训练领域之间的对应关系。同时，需要识别同一文本上不同质量指标评价明显不一致的情况，定义质量冲突并分析其产生原因，在抽样集和扩展集上检验冲突规律的稳定性。在此基础上，利用 $A_4$--$A_{15}$ 中的 17 维领域配比和交叉熵损失数据，建立领域配比向量 $\boldsymbol{p}$ 与 Loss 之间的定量关系，其中
\begin{equation}
  \boldsymbol{p}=(p_1,p_2,\ldots,p_{17}),\qquad
  p_i\geq 0,\qquad \sum_{i=1}^{17}p_i=1.
  \label{eq:problem-restatement-simplex}
\end{equation}
需要分析不同领域及其组合对模型性能的影响，并通过检验数据和外推数据验证模型的可靠性与稳健性。

\subsubsection{问题二：跨维度数据融合与广义标度律}

问题二要求在经典标度律的基础上，将问题一得到的数据质量 $Q$ 和领域配比 $\boldsymbol{p}$ 纳入模型，建立可以同时考虑参数量 $N$、训练数据量 $D$、数据质量和领域组成的广义标度律。模型拟合主要使用附件 $B_1$ 中的真实训练轨迹，并结合 $B_2$、$B_3$ 进行轨迹或模型外部验证，利用 $B_4$、$B_5$ 检验模型在跨模型和文献数据上的泛化能力，再通过 $B_6$--$B_8$ 的质量补充实验分析数据质量的作用，利用 $B_9$、$B_{10}$ 讨论大规模模型的外推边界。在模型估计和验证的基础上，还需要比较参数规模、数据规模、数据质量和领域配比的边际效用与弹性，分析不同领域之间的替代或互补关系，并推导质量提升与参数规模扩张之间的可计算等价条件，同时明确半合成数据、插值数据和外推数据的可信度边界。

\subsubsection{问题三：算力约束下的多维资源联合优化与结构性转移}

问题三研究有限算力预算下的资源配置问题，需要在参数量 $N$、训练数据量 $D$、数据质量 $Q$ 和领域配比 $\boldsymbol{p}$ 之间进行联合分配，使模型预测 Loss 尽可能低。优化模型应同时考虑基础训练成本、数据质量提升成本和长上下文注意力成本，并比较指数型、幂函数型和对数渐进型质量成本函数对最优解的影响。上下文长度 $L_{\mathrm{ctx}}$ 由模型架构和任务需求外生给定，需要依据附件 $C_7$ 确定其合理范围并进行敏感性分析，同时设置低、中、高不同量级的算力预算，考察预算变化时各类资源的最优分配比例是否发生明显改变，进一步给出结构性转移的判定标准以及训练成本与注意力成本相当的临界条件。

\subsubsection{问题四：技术演进分析与模型能力前沿预测}

问题四利用附件 $C$ 中的模型规模、训练算力、训练数据量、发布时间、开源状态和多维 Benchmark 数据，分析大语言模型能力提升的来源，并预测未来能力的前沿边界。首先，需要将能力提升分解为参数规模、数据规模和算力增长带来的规模扩张效应，以及架构改进、数据工程和对齐方法等带来的非规模技术进步，并分别量化和比较两类因素的贡献。其次，需要明确综合 Benchmark 指标、开源模型筛选标准、\emph{pretrained} 与 \emph{chat/finetuned} 模型的分类方式以及时间轴口径，并利用 $C_5$ 或 $C_6$ 建立 Loss 与 Benchmark 之间的桥接关系，分析桥接误差对结论的影响。最后，结合 $C_1$--$C_4$ 的总体信息和 $C_8$ 的逐任务评测结果，在算力增长放缓的情景下预测未来 12 个月或 24 个月开源大语言模型的能力前沿边界，并给出相应的不确定性分析。
 引入。
% 仅包含问题背景与问题提出，不包含具体模型、计算结果或结论。

\section{问题重述}
\label{sec:problem-restatement}

\subsection{问题背景}

近年来，大语言模型在自然语言处理、科学研究和智能应用等领域快速发展，经证明发现模型能力通常受到参数规模、训练数据量和训练算力的共同影响。其中，标度律是用以刻画模型损失随规模变化的经验关系，而经典标度律则是主要用于描述参数量 $N$、训练数据量 $D$ 与验证集交叉熵损失 Loss 之间的关系。但是随着大语言模型的长期实践发现，模型性能并不只取决于规模因素，训练数据质量、不同领域数据的配比、上下文长度以及模型架构等因素同样会产生重要影响。特别是在高质量语料日益稀缺、数据污染和“数据墙”问题逐渐突出的背景下，仅依靠扩大模型规模或增加训练数据，已经难以全面解释模型能力的变化。

但是与此同时，增加参数量、扩充训练数据、提升数据质量以及处理更长上下文都会带来额外的算力消耗；故而在算力预算有限的情况下，模型规模、数据规模、数据质量和领域配比之间存在资源竞争与相互替代关系。因此，本题要求综合利用附件中的数据质量信号、配方实验、模型训练日志、历史模型元数据和多维 Benchmark 评测数据，建立“数据质量与领域配比—模型 Loss—模型能力—算力配置”之间的数学联系。四个问题依次研究数据质量和领域配比、广义标度律、算力约束下的资源优化以及模型能力演进预测，前一问题的分析结果将作为后一问题的重要基础。

\subsection{问题提出}

\subsubsection{问题一：数据质量评价、质量冲突消解与领域配比建模}

问题一主要研究训练数据质量及领域配比对模型性能的影响。利用附件 $A_1$--$A_3$ 中的 22 个质量指标，需要对不同方向的指标进行预处理、标准化和正负向统一，构造样本级、语料级或领域级质量评分 $Q$，并依据 $A_{16}$ 建立质量信号领域与 17 个训练领域之间的对应关系。同时，需要识别同一文本上不同质量指标评价明显不一致的情况，定义质量冲突并分析其产生原因，在抽样集和扩展集上检验冲突规律的稳定性。在此基础上，利用 $A_4$--$A_{15}$ 中的 17 维领域配比和交叉熵损失数据，建立领域配比向量 $\boldsymbol{p}$ 与 Loss 之间的定量关系，其中
\begin{equation}
  \boldsymbol{p}=(p_1,p_2,\ldots,p_{17}),\qquad
  p_i\geq 0,\qquad \sum_{i=1}^{17}p_i=1.
  \label{eq:problem-restatement-simplex}
\end{equation}
需要分析不同领域及其组合对模型性能的影响，并通过检验数据和外推数据验证模型的可靠性与稳健性。

\subsubsection{问题二：跨维度数据融合与广义标度律}

问题二要求在经典标度律的基础上，将问题一得到的数据质量 $Q$ 和领域配比 $\boldsymbol{p}$ 纳入模型，建立可以同时考虑参数量 $N$、训练数据量 $D$、数据质量和领域组成的广义标度律。模型拟合主要使用附件 $B_1$ 中的真实训练轨迹，并结合 $B_2$、$B_3$ 进行轨迹或模型外部验证，利用 $B_4$、$B_5$ 检验模型在跨模型和文献数据上的泛化能力，再通过 $B_6$--$B_8$ 的质量补充实验分析数据质量的作用，利用 $B_9$、$B_{10}$ 讨论大规模模型的外推边界。在模型估计和验证的基础上，还需要比较参数规模、数据规模、数据质量和领域配比的边际效用与弹性，分析不同领域之间的替代或互补关系，并推导质量提升与参数规模扩张之间的可计算等价条件，同时明确半合成数据、插值数据和外推数据的可信度边界。

\subsubsection{问题三：算力约束下的多维资源联合优化与结构性转移}

问题三研究有限算力预算下的资源配置问题，需要在参数量 $N$、训练数据量 $D$、数据质量 $Q$ 和领域配比 $\boldsymbol{p}$ 之间进行联合分配，使模型预测 Loss 尽可能低。优化模型应同时考虑基础训练成本、数据质量提升成本和长上下文注意力成本，并比较指数型、幂函数型和对数渐进型质量成本函数对最优解的影响。上下文长度 $L_{\mathrm{ctx}}$ 由模型架构和任务需求外生给定，需要依据附件 $C_7$ 确定其合理范围并进行敏感性分析，同时设置低、中、高不同量级的算力预算，考察预算变化时各类资源的最优分配比例是否发生明显改变，进一步给出结构性转移的判定标准以及训练成本与注意力成本相当的临界条件。

\subsubsection{问题四：技术演进分析与模型能力前沿预测}

问题四利用附件 $C$ 中的模型规模、训练算力、训练数据量、发布时间、开源状态和多维 Benchmark 数据，分析大语言模型能力提升的来源，并预测未来能力的前沿边界。首先，需要将能力提升分解为参数规模、数据规模和算力增长带来的规模扩张效应，以及架构改进、数据工程和对齐方法等带来的非规模技术进步，并分别量化和比较两类因素的贡献。其次，需要明确综合 Benchmark 指标、开源模型筛选标准、\emph{pretrained} 与 \emph{chat/finetuned} 模型的分类方式以及时间轴口径，并利用 $C_5$ 或 $C_6$ 建立 Loss 与 Benchmark 之间的桥接关系，分析桥接误差对结论的影响。最后，结合 $C_1$--$C_4$ 的总体信息和 $C_8$ 的逐任务评测结果，在算力增长放缓的情景下预测未来 12 个月或 24 个月开源大语言模型的能力前沿边界，并给出相应的不确定性分析。
